\section{Protocol and Factorial Design}
\label{sec:protocol}

\begin{quote}
\small\raggedright
Full machine-readable protocol: \texttt{RESEARCH\_PROTOCOL.md} v6.5 and
\texttt{STAGE\_CONTRACTS.md}; this section is the paper-facing summary.
\end{quote}

\subsection{Research questions and gates}
\begin{description}
\item[RQ1 (attribution).] Which framing component(s) causally increase attack
success in \LALMs{}, and are effects robust across language, model, and
out-of-distribution inputs?
\item[RQ2 (cross-modal).] Do text and audio renderings of the same input
diverge in safety judgment, and is the divergence model-dependent?
\item[RQ3 (detection).] Can a structured, signal-level fusion framework
discriminate framed attacks in-distribution, and how does it behave under
adaptive attack and audio degradation?
\end{description}

We pre-registered two decision gates. \textbf{G1} (attribution established)
requires a consistent, $\geq$10~pp effect whose bootstrap CI excludes zero under
the dual-judge caliber with acceptable scoring dispute. \textbf{G2} (defense
value) requires a significant TPR gain, $\geq$3 contributing branches, and a
benign false-positive rate $\leq$5\%. As reported in
Sections~\ref{sec:results} and \ref{sec:discussion}, G1 passes and
\textbf{G2 does not pass}; the paper therefore presents the attribution and
mechanism results as primary and the fusion framework as an application-level
validation with explicit boundary conditions.

\subsection{Factors and manipulation}
A paired factorial design \citep{fisher1935} varies four framing components
while holding the harmful request fixed:

\begin{table}[t]
\centering
\small
\caption{\textbf{The four framing factors.} Varied in the paired factorial
design (\S\ref{sec:setup}).}
\label{tab:factors}
\begin{tabular}{@{}l@{\hspace{2em}}p{0.30\textwidth}@{\hspace{2em}}p{0.28\textwidth}@{}}
\toprule
Factor & Definition & Levels \\
\midrule
\Et{} & Narrative text framing (emotional / narrative wording)
      & neutral vs.\ narrative \\
\Nfac{} & Narrative structure (event-chain packaging)
        & direct request vs.\ chained narrative \\
\Rfac{} & Role assignment (identity + task)
        & none vs.\ expert role \\
\Asfac{} & Acoustic style (TTS rendering)
         & neutral voice vs.\ styled voice \\
\bottomrule
\end{tabular}
\end{table}

Three templates instantiate each condition (9 templates total); a manipulation
check on response-level markers confirms each textual factor was delivered
(mean marker deltas 0.37--0.60, threshold 0.95 keep rate 1.0). We disclose that
the acoustic factor is implemented via TTS voices (a female neutral voice vs.\ a
male styled voice), so it confounds speaker identity and gender with prosody.
We do not claim \Asfac{} isolates prosody, and we verify no systematic length
difference across \Asfac{} variants (mean difference 1.18 chars, $p=0.24$).
